\begin{figure*}[t]
\centering
\includegraphics[width=0.49\textwidth]{fig6_fullq_positive.pdf}\hfill
\includegraphics[width=0.49\textwidth]{fig6_fullq_negative.pdf}
\caption{Previously calculated complete connected spectra for $L=40$ open
chains, retained without changing the two selected panels.
$J_2/J_1=0.5$ (left), $-0.5$ (right), and $J_3/J_1=0.5$.
The triangular time window at $TJ_1=60$ has width $0.09277J_1$;
the production time step and bond dimension are $\Delta tJ_1=0.05$ and 192.
Both use the same display coordinate
$u(s)=\operatorname{asinh}(s/0.02)/\operatorname{asinh}(6/0.02)$ for
$s=J_1S_T^{zz}\in[0,6]$. This monotonic display does not change the spectral
arrays or their resolution. No Gaussian smoothing, momentum interpolation,
or reflection averaging is applied. The 41 momenta were independently
calculated. These finite-coupling maps are contextual, not quantitative
evidence for the weak-coupling retuned residual.}
\label{fig:rank-original}
\end{figure*}

\section{Local root and numerical protocols}

For the quartet root at $q=\pi$, the highest-weight $M=3/2$ state has
single-down-spin amplitudes $(u,v,w,v,u)$, with $w=-2(u+v)$ from the
highest-weight condition. The zero of the projected alternating probe gives
$v^2=1/8$ after normalization. The relevant eigenvector ratio is
$x=u/v=(\sqrt7-2)/3$, with
$\alpha=-x(3+2x)/(1-x^2)$, giving Eq.~\eqref{eq:rank-exact-root}.
All local gaps and resolvents were also obtained from direct diagonalization
of the 32-dimensional pentamer Hamiltonian. No reduced transition amplitudes
were fitted to the microscopic spectra.

Microscopic Hamiltonians were built in fixed-$S^z$ bases. The three-cell
comparison used a real symmetric sparse eigensolver with tolerance
$10^{-13}$, at least 50 Krylov vectors, and four extra eigenpairs above the
expected retained band. The numerical acceptance tests required two ground
states in the chosen sector, a complete retained band, and eigen-residuals
below $10^{-10}$. The anticipated small-coupling slope intervals were
$[1.8,2.2]$ and $[3.7,4.3]$; the quartet's weakest-coupling coefficient was
required to agree within $5\%$. These are deterministic validation tests,
not uncertainty intervals. All passed. The parameter and momentum search
used to choose the roots was exploratory; the subsequent microscopic tests
used those fixed predictions.
The restart check used tolerance $10^{-14}$, 120 Krylov vectors, a different
seed, and six extra eigenpairs above the retained count at $g=0.005,0.01$.
It used the same Hamiltonian assembly and is a numerical-sensitivity check,
not an independent computational method.

For the four-cell quartet test, the $S^z=0$ microscopic basis has 184756
states. We obtained 45 eigenpairs, including all 44 retained-sector states
(the ground state and 43 excited states), with tolerance $10^{-13}$.
The two smallest couplings used 140 Krylov vectors; a repeat at
$g=0.000625$ used 180 and a different random starting vector. The line-vector
error compares weights in increasing-energy order, including a zero
ground-state entry, after division by $\beta^2$. It is not a test of a
broadened line shape. The strongest-coupling points were diagnostic tests,
not part of the weak-domain acceptance claim.

Effective-chain ED used the $S^z=0$ sector, whose dimension reaches 116304
for ten spin-$3/2$ sites. The norm and overlap in
Eq.~\eqref{eq:rank-minimum} were evaluated directly on the ground-state
sources, without summing an incomplete set of excited states. The lowest
line and the energy moment were checked separately. The periodic scalar
source and the sum of the open-chain decomposition supplied internal
operator checks.

The longer open effective chains were computed by DMRG~\cite{white1992,schollwoeck2011}
using ITensors/ITensorMPS~\cite{fishman2022}
with conserved $S^z$, explicit spin-$3/2$ local operators, and MPO contractions
of $D^\dagger D$, $K^\dagger K$, and $D^\dagger K$. Thus there is no additional
truncation from constructing an applied source MPS. Refinement used bond
dimensions $128,192,256$, with 24, 16, and 16 sweeps, cutoff $10^{-13}$,
and eigensolver tolerance $10^{-13}$. Larger-bond runs started from the
previous state without decreasing its bond dimension. Acceptance required
relative change of $\Cmin$ below $1\%$, energy change per cell below
$10^{-8}$, energy variance per cell below $10^{-8}$, and maximum local
$|\langle T_n^z\rangle|<10^{-4}$. The final $N=20$ variance was
$1.21\times10^{-9}$ and its maximum local moment was $1.17\times10^{-9}$.
The earlier unconverged runs and the failed $N=40$ refinement are retained
with their original criteria.

Periodic DMRG used the same operators with the closing bond included in the
MPO. The $N=6$ source norm was cross-checked against ED. Bond dimensions
$128,256,384$ were followed by 512 for $N=12$, with cutoff $10^{-13}$ and
16 warm-start sweeps at each refinement. In addition to the preceding
energy and weight criteria, acceptance required
$|\langle\bm T_{\mathrm{tot}}^2\rangle|<10^{-4}$, bond-energy spread below
$10^{-5}$, and scalar-removal and energy-reconstruction identities within
$10^{-10}$. The first $N=12$ batch failed the energy-change test at 384;
the subsequent 512 calculation passed without changing that threshold.
All sizes here count effective spin-$3/2$ sites, not microscopic spins.
The bounded refinement also tested $N=16$ at dimensions $256,512,768$;
the last run failed the unchanged energy-change threshold and no larger
periodic calculation was started.

For the $Q=2\pi/3$ quartet check, the same periodic ground states can be
reused: changing the local root multiplies the leading uniform Heisenberg
Hamiltonian by the positive scalar $c_1(\alpha)^2$ without changing its
eigenvectors. The probe coefficients were recomputed at the new root,
including their complex Fourier phases. Connected norms use
$\langle D^\dagger K\rangle-\langle D\rangle^*\langle K\rangle$,
and the exchange retuning remains real. The six-cell MPO result reproduced
ED to relative error $4.5\times10^{-15}$. The twelve-cell weight passed the
$1\%$ bond-dimension-change test, and the imaginary part of the normalized
source overlap was below $10^{-8}$. This is an effective-model measurement,
not microscopic sixty-spin diagonalization.

\section{Thermal traces and exact local checks}
\label{app:exact-local-check}

The four- and six-cell thermal effective calculations included every
$S^z$ sector, with total dimensions $4^N$ or $2^N$. For six quartet
cells the largest block had dimension 580. Real and imaginary operator
components were transformed into the complete energy bases, and all
transitions with $|\epsilon_f-\epsilon_i|<10^{-9}$ were treated as
equal-energy transitions. Checks included Hermiticity, trace sum rules,
eigen-residuals below $10^{-11}$, detailed balance, and agreement between
direct residual norms and the thermal least-squares expression. A separate
256-dimensional tensor-product assembly reproduced the four-cell result.
Random unitary changes within each degenerate eigenspace left the thermal
weights invariant. Both floating-point assemblies used the same dense
diagonalization library; they are not independent solution methods.

\paragraph*{Exact microscopic prefactors.}
At the quartet root, an unnormalized highest-weight cell state in the
one-down-spin basis, ordered by the position of that spin, is
\begin{equation}
 v=(1,\ 2+\sqrt7,\ -6-2\sqrt7,\ 2+\sqrt7,\ 1)^{\mathsf T}.
\end{equation}
Its squared norm is $88+32\sqrt7$ and its energy is
$E_5=-(5+4\sqrt7)/12$. Total-spin lowering constructs the quartet
projector $P_{\rm cell}$. We evaluated the reduced resolvent
$G=(H_{\rm cell}-E_5+P_{\rm cell})^{-1}-P_{\rm cell}$ separately in all
six magnetization blocks, using exact arithmetic in $\mathbb Q(\sqrt7)$.
An exact $LDL^{\mathsf T}$ factorization of
$H_{\rm cell}-E_5-\tfrac18(1-P_{\rm cell})+P_{\rm cell}$
has strictly positive diagonal pivots in every block. This certifies
that the quartet is the ground multiplet and that $\Dloc>1/8$.
The projected tensor in Eq.~\eqref{eq:rank-local-tensor} is
$s_0\delta_{z\mu}I+t_0Q^{z\mu}$, with
\begin{align}
 s_0&=\frac{161590\sqrt7-361001}{180792},\notag\\
 t_0&=\frac{169277-70390\sqrt7}{45198}.
\end{align}
Both endpoint choices and all $z,+,-$ components satisfy this identity
exactly. Multiplication by $-c_1$ gives $a,b$ in
Eq.~\eqref{eq:exact-local-coefficients}. The elementary bounds
$2.645<\sqrt7<2.646$ certify $b>20086103/162712800>0$.
Comparison with the separately assembled floating-point two-cell
resolvent gives an absolute coefficient difference below
$3.0\times10^{-16}$. This is a check of the same perturbative formula,
not a distinct many-body approximation.

The other exact checks used rational arithmetic in the unnormalized local basis
$|k\rangle=(T^-)^k|s,s\rangle$, where
$T^+|k\rangle=k(2s-k+1)|k-1\rangle$ and
$\langle k|k\rangle=k!(2s)!/(2s-k)!$.
This gives exact zero tests for the partial traces in
Eq.~\eqref{eq:thermal-bound}, with no numerical threshold.
For the four-cell activation calculation, the two sublattice spins were
coupled from $j_A=j_B=3$. In an unnormalized product basis, a highest-weight
total-spin-$J$ state has coefficients $c_k$ at sublattice lowering degrees
$(k,6-J-k)$, obeying
\begin{equation}
 \frac{c_{k+1}}{c_k}
 =-\frac{(6-J-k)(J+1+k)}{(k+1)(6-k)}.
\end{equation}
This recurrence, followed by total-spin lowering, constructs the ground
singlet and first triplet without a floating-point eigensolver.
Acting with $Y+3\widetilde D$ annihilates the singlet. On the triplet,
the residual separates into final eigenvalues $-9$ and $-4$ of $h$;
projection by $\bm A^2$ resolves the two $-4$ channels. Their exact norms
give Eq.~\eqref{eq:thermal-activation}.

The ten-spin microscopic thermal calculation used all 11 $S^z$ blocks,
whose largest dimension is 252. The full partition function was retained.
Both initial and final states in $W_+$ were restricted to
$E-E_0<\Dloc/2$; there were 16 retained quartet or 4 retained doublet
states, as required. Eigen-residuals were below $3.4\times10^{-15}$,
and omitted-band thermal populations were below $1.1\times10^{-57}$
for the tested parameters. The temperature was
$T=\tau\beta c_1(\alpha_*)^2$, including on the retuned trajectory.

For the finite momentum interval we formed the low-band matrix elements
of each microscopic $S_j^z$. At each quadrature point the complex Fourier
amplitudes were combined before taking their absolute squares; this
avoids subtracting nearly equal contributions to a tiny central intensity.
Uniform-window integrations with 16 and 32 Gauss--Legendre nodes agreed
within $10^{-7}$ relatively. Zero-width values reproduced the separate
microscopic calculation within $1.2\times10^{-10}$. This is integration
over a specified momentum acceptance, not Gaussian broadening or an
image-processing operation.

For twenty microscopic spins, the sparse $S^z=0$ Hamiltonian has dimension
184756. We computed 48 lowest eigenpairs and verified that precisely 44
lie below $E_0+\Dloc/2$. Within this low band, diagonalization of
$(S^+V)^\dagger(S^+V)$ resolves total spin. Rediagonalizing the Hamiltonian
within each spin sector and applying normalized total-spin raising gives
all positive-$S^z$ states; spin reversal supplies the negative sectors.
The multiplicities of $J=0,\ldots,6$ are $4,9,11,10,6,3,1$,
giving 256 states. Every raised sector was checked against an independently
assembled microscopic Hamiltonian. Residuals were below $10^{-12}$ and
orthogonality errors below $10^{-11}$. The same procedure at ten spins
reproduced the full-spectrum thermal weight to relative error below
$10^{-7}$ before the twenty-spin run.
At the weakest coupling, an independent restart with a different random
vector, a larger Krylov space (220 instead of 180), and tighter tolerance
changed the total, ground-initial, and excited-initial weights by at most
$8.0\times10^{-8}$ relatively across the tested temperatures and both
frequency definitions. This is solver sensitivity testing, not a distinct
many-body method.

The omitted states have a rigorous thermal bound. Relative to uncoupled
cells at $\alpha_*$, the perturbation norm is at most
$v=(3/4)[N|\beta|+2N|\alpha-\alpha_*|]$.
The verified condition $4v<\Dloc$ isolates the retained band and places
it below the chosen cutoff. Weyl bounds put every omitted state at least
$\Dloc-2v$ above the true ground energy. With the retained partition
function shifted by that energy,
\begin{equation}
 \frac{Z_{\rm out}}{Z_{\rm low}}
 \le \frac{(2^{5N}-4^N)e^{-(\Dloc-2v)/T}}{Z_{\rm low}}.
\end{equation}
The logarithm of this bound is below $-2167$ in all tested thermal cases;
the tail was bounded, not set to zero by numerical underflow.
The raw positive-frequency cutoff was $10^{-10}$ in $J_1$ units.
An audit at $10^{-11}$ and $10^{-9}$ exposed the quasielastic contribution
at $\tau=1.6$. The subsequent fixed scaled windows are reported separately
and do not replace the initial protocol or its data.

\section{Relation to the original full spectra}

The original $L=40$ maps in Fig.~\ref{fig:rank-original} have
$\beta/\Dloc\simeq1.46$ and $3.96$, outside the validated perturbative
range, and finite-time resolution broader than the separately calculated
open-chain gaps. They neither establish the weak-coupling powers in
Fig.~\ref{fig:rank-evidence} nor resolve the tiny retuned weights; they
are retained as context, not validation of the composite-probe expansion.

